% !TeX root = ../main.tex

\chapter{Összegzés és jövőbeli fejlesztések}

\section{Az elért eredmények értékelése}
A Smart Turntable Assistant sikeresen eléri elsődleges célját, az analóg bakelitlemez-hallgatási élmény modernizálását. A nyers elektromos audiojel passzív monitorozásával a rendszer áthidalja a fizikai média és a digitális ökoszisztéma közötti szakadékot. Az alkalmazás közel 80\% pontossággal azonosítja a lejátszott sávokat, a hallgatási munkamenetet nagy felbontású borítóképekkel és szinkronizált dalszövegekkel gazdagítja, és zökkenőmentesen naplózza a hallgatási előzményeket a külső közösségi zenei platformokra.

Továbbá, az audiofil szintű diagnosztikai eszközök beépítése valós idejű vizuális visszajelzést biztosít a felhasználóknak a tű és a lemezek fizikai állapotáról. A rendkívül moduláris, szétválasztott (decoupled) architektúra biztosítja, hogy a mögöttes Python backend és a reaktív webes interfész hatékonyan működjön egy olyan professzionális és robusztus szoftveres megoldásban, amely tiszteletben tartja a hagyományos hanglemezhallgatás élményét, miközben modern digitális kényelmet nyújt.

\section{Jövőbeli fejlesztési lehetőségek}
Bár a jelenlegi implementáció teljesíti az összes kezdeti követelményt, számos lehetőség kínálkozik a jövőbeli továbbfejlesztésre. A legjelentősebb fejlődés saját neuron hálók használatával való sibilance detection és két zeneszám közti különbség (tempó és hangszín) detektálásával lenne elérhető.

Ezenfelül a diagnosztikai csomag kibővíthető lenne egy automatizált lemezjátszó-kalibrációs móddal is. Egy lemez lejátszásával ki tudja választani a felhasználó, hogy hol megy már a zene, ezzel a music detection algoritmus érzékenységét egyből be tudja állítani egy elfogadható állapotra, minimalizálva a manuális konfigurációt.